\begin{proposition}[6.3.8]
\label{IV.2.s6.3.8-fr}
Sea $A$ un anillo local noetheriano; supongamos que existe un $A$-módulo $M$
de tipo finito que sea un $A$-módulo de Cohen--Macaulay y cuyo soporte sea
igual a $\operatorname{Spec}(A)$. Sea $B$ un anillo cociente de $A$, y sea
$f:\operatorname{Spec}(\widehat B)\to\operatorname{Spec}(B)$ el morfismo
canónico; entonces, para todo $x\in\operatorname{Spec}(B)$, $f^{-1}(x)$ es un
esquema de Cohen--Macaulay.
\end{proposition}

Como $B$ es un $A$-módulo de tipo finito, se tiene
$\widehat B=B\otimes_A\widehat A$ $(\mathbf{0}_{\mathrm I}, 7.3.3)$, luego $f$ es la
restricción a $\operatorname{Spec}(\widehat B)$ del morfismo canónico
$\operatorname{Spec}(\widehat A)\to\operatorname{Spec}(A)$, y las fibras de
estos dos morfismos en un punto de $\operatorname{Spec}(B)$ son, por tanto,
las mismas. La cuestión se reduce así a demostrar la proposición cuando $B=A$.
Sea, pues, $\mathfrak p=\mathfrak j_x$ un ideal primo de $A$; como, por
hipótesis, $\mathfrak p\in\operatorname{Supp}(M)$, se sabe (\textbf 0, 16.5.6 y
16.5.9) que existe una sucesión $M$-regular $(t_i)_{1\leq i\leq r}$
de elementos de $\mathfrak p$ tal que
$N=M/(\sum_{i=1}^{r}t_iM)$
sea un $A$-módulo de Cohen--Macaulay, $\mathfrak p$ un ideal primo asociado
\emph{minimal} de $N$ y
$\dim(N)=\dim(M/\mathfrak pM)=\dim(A/\mathfrak p)$. El mismo razonamiento
que al principio muestra que se puede sustituir $M$ por $N$ y $A$ por
$A/(\sum_{i=1}^{r}t_iA)$ y, por tanto, suponer que $\mathfrak p$ es un

\oldpage[IV-2]{141}

ideal primo \emph{minimal} de $A$. Pongamos, para simplificar,
$A'=\widehat A$, $M'=\widehat M=M\otimes_AA'$ $(\mathbf{0}_{\mathrm I}, 7.3.3)$; se
sabe (\textbf 0, 16.5.2) que $M'$ es un $A'$-módulo de Cohen--Macaulay, lo que implica
que, para todo ideal primo $\mathfrak p'$ de $A'$,
$M'_{\mathfrak p'}$ es un $A'_{\mathfrak p'}$-módulo de Cohen--Macaulay
(\textbf 0, 16.5.10). Teniendo en cuenta (\textbf I, 3.6.5), se ve que, si se pone
$A''=A'\otimes_AA_{\mathfrak p}$ y
$M''=M'\otimes_AA_{\mathfrak p}$, $M''$ es un $A''$-módulo de
Cohen--Macaulay. Para todo ideal primo $\mathfrak p''$ de $A''$ situado sobre
$\mathfrak p$,
$M''_{\mathfrak p''}=M_{\mathfrak p}\otimes_{A_{\mathfrak p}}
A''_{\mathfrak p''}$ es, por tanto, un $A''_{\mathfrak p''}$-módulo de
Cohen--Macaulay (\textbf 0, 16.5.10); por otra parte, $M_{\mathfrak p}$ es, por
hipótesis, un $A_{\mathfrak p}$-módulo de Cohen--Macaulay y
$A''_{\mathfrak p''}$ es un $A_{\mathfrak p}$-módulo \emph{plano}, puesto que $A'$
es un $A$-módulo plano $(\mathbf{0}_{\mathrm I}, 7.3.3\text{ y }6.3.2)$. Se concluye
de (6.3.3) que, si $k$ es el cuerpo residual de $A_{\mathfrak p}$,
$k\otimes_{A_{\mathfrak p}}A''_{\mathfrak p''}$ es un anillo de
Cohen--Macaulay. Pero, puesto que $A_{\mathfrak p}$ es de dimensión $0$, los
ideales primos de $A''$ corresponden biunívocamente a los de
$k\otimes_{A_{\mathfrak p}}A''$ (\textbf I, 3.5.7), y si $\mathfrak q''$ es el ideal
primo de $k\otimes_{A_{\mathfrak p}}A''$ correspondiente a
$\mathfrak p''$, los anillos locales
$(k\otimes_{A_{\mathfrak p}}A'')_{\mathfrak q''}$ y
$k\otimes_{A_{\mathfrak p}}A''_{\mathfrak p''}$ son isomorfos. Por tanto
(\textbf 0, 16.5.13), el anillo $k\otimes_{A_{\mathfrak p}}A''$ es un anillo de
Cohen--Macaulay. C.Q.F.D.
